\bookchapter{第四十七章：速查手冊與資源}

\chapterintro{本章是一份可以獨立使用的速查手冊。當你記不清某個原則或需要快速決策時，直接翻到這裡。建議你把它當成“口袋卡片”：五大原則速查、工具決策樹、常見問題、學習資源，一頁不落。}

\booksection{一、思維模型速查卡}

\begin{center}
\begin{tabularx}{\linewidth}{|X|X|X|}
\hline
\textbf{原則} & \textbf{一句話} & \textbf{操作要點} \tabularnewline
\hline
Think Before Coding & 先想清楚再開口 & 回答三問：做什麼？結果長啥樣？有什麼約束？ \tabularnewline
Simplicity First & 用最少程式碼解決問題 & 每次只做一個功能，砍掉不必要的複雜度 \tabularnewline
Surgical Changes & 精準修改不越界 & 明確指定修改範圍，拒絕AI的"順手改一下" \tabularnewline
Goal-Driven Execution & 說"做什麼"不說"怎麼做" & 給目標不給步驟，讓AI自主選擇路徑 \tabularnewline
Iterate \& Verify & 生成-驗證-修正迴圈 & 三層驗證：功能、邊界、架構 \tabularnewline
\hline
\end{tabularx}
\end{center}

\begin{center}\small\color{gray}五大思維原則速查卡\end{center}

\booksection{二、工具選擇決策樹}

當你需要選擇 AI 程式設計工具時，按照以下決策路徑選擇。注意：這不是“哪個最好”的排名，而是“哪個最適合你”的判斷框架。

你要做什麼？

零基礎，想快速體驗 選雲端 Builder (v0/Bolt/Replit)

日常開發，需要程式碼補全 選 IDE 內建 AI (Cursor/Trae)

大型專案，需要批次操作 選終端 Agent (Claude Code)

快速原型驗證想法 選雲端 Builder 或 Trae SOLO

\begin{center}\small\color{gray}圖 6：AI 程式設計工具選擇決策樹\end{center}

\booksection{三、常見問題 FAQ}

Q1：我完全不會程式設計，能用 AI 程式設計嗎？

能。2026 年的 AI 程式設計工具已經足夠強大，你完全可以用自然語言描述需求，讓 AI 生成程式碼。但有兩點建議：第一，從雲端 Builder 工具（如 v0、Bolt）開始，它們不需要配置開發環境；第二，邊用邊學基礎概念——不需要精通語法，但你需要能“看懂”AI 寫的程式碼的大意，這樣才能有效驗收。

Q2：我應該學哪個程式語言？

取決於你要做什麼。做網站前端——JavaScript/TypeScript；做資料處理或 AI——Python；做企業級後端——Java 或 Go。但更重要的是：在 AI 程式設計時代，“語言”的重要性在下降。AI 能幫你處理語法細節，你需要理解的是“程式的結構”——什麼是函式、什麼是迴圈、什麼是資料結構。這些概念是跨語言的，一旦理解，換什麼語言都快。

Q3：AI 寫的程式碼有版權問題嗎？

截至 2026 年 8 月，這仍是一個有爭議的法律問題。一般來說，AI 生成的程式碼在大多數司法管轄區不被視為有明確的版權所有者。但如果你用於商業專案，建議：(1) 瞭解你使用的 AI 工具的服務條款中關於智慧財產權的條款；(2) 對關鍵程式碼進行人工審查和修改，使其具有你的創造性投入；(3) 避免讓 AI 直接複製可能有版權的開原始碼。

Q4：AI 會取代程式設計師嗎？

短期內不會。AI 取代的不是“程式設計師”這個職業，而是“逐行編碼”這項工作。就像計算機沒有取代數學家，而是讓數學家從繁瑣的手算中解放出來一樣。AI 程式設計工具讓程式設計師從“寫程式碼”轉向“想問題”——這實際上提高了程式設計師的價值上限。未來的競爭力不在於“能寫多少行程式碼”，而在於“能解決多複雜的問題”。

Q5：如何判斷 AI 的產出是否正確？

三層驗證法。第一層：功能驗證——它能不能跑？輸入輸出對不對？這是最基本的。第二層：邊界驗證——空輸入、超長輸入、異常輸入會怎樣？這能發現隱藏的問題。第三層：架構驗證——程式碼結構是否合理？有沒有硬編碼？是否便於擴充套件？這需要你有一定的技術判斷力。如果你無法判斷第三層，至少做到前兩層——能跑、邊界正常，就比大多數初學者強了。

\booksection{四、推薦學習資源}

以下是截至 2026 年 8 月的推薦學習資源。請注意，網際網路資源連結可能變化，如果某個連結失效，搜尋其標題通常能找到最新版本。

\begin{center}
\begin{tabularx}{\linewidth}{|X|X|X|}
\hline
\textbf{型別} & \textbf{資源} & \textbf{推薦理由} \tabularnewline
\hline
原始文獻 & Andrej Karpathy 的 CLAUDE.md & 五大原則的源頭，70行但字字珠璣 \tabularnewline
原始文獻 & Anthropic 官方檔案 (Claude Code) & 終端Agent的最佳實踐指南 \tabularnewline
教學 & Vibe Coding 社群整理的 Skills 合集 & 社群實踐者總結的實戰技巧集合 \tabularnewline
教學 & Cursor 官方檔案和 YouTube 頻道 & IDE內建AI的權威使用指南 \tabularnewline
社群 & GitHub 上的 awesome-claude-code & Claude Code 的最佳實踐和外掛集 \tabularnewline
社群 & Reddit 的 r/cursor 社群 & 全球Cursor使用者的經驗分享和討論 \tabularnewline
基礎 & MDN Web Docs (developer.mozilla.org) & 前端開發最權威的參考檔案 \tabularnewline
基礎 & Python 官方教學 (docs.python.org) & Python學習的官方起點 \tabularnewline
\hline
\end{tabularx}
\end{center}

\begin{center}\small\color{gray}推薦學習資源清單\end{center}

\booksection{五、結語：工具在手，方法論在心}

回到本書的核心信念：工具每三個月迭代一次，方法論每三年才進化一次。在 2026 年 8 月的今天，AI 程式設計工具的能力已經強大到令人驚歎——你能用一句話建立一個網站，用一段描述生成一個全棧應用，用一條指令完成一個複雜重構。但工具越強大，方法論的重要性反而越突出——因為工具的強大放大了你的每一個決策。

如果你的需求清晰，AI 能幫你高效實現；如果你的需求模糊，AI 也會高效地做出錯誤的東西。如果你遵循簡單至上，AI 能幫你快速驗證想法；如果你貪多求全，AI 也會幫你快速堆積技術債務。如果你精準修改，AI 能幫你外科手術般地最佳化程式碼；如果你指令模糊，AI 也會幫你大刀闊斧地引入 bug。

這就是為什麼本書花了大部分篇幅講方法論而非工具操作。不是因為工具不重要，而是因為方法論是工具的“乘數”——方法論越好，工具的價值越大；方法論為零，工具再強也只是零乘以任何數。

你不再是一個人在程式設計。你是一個指揮官，AI 是你的軍團。你的思考質量決定了軍團的戰鬥力。從今天開始，用方法論武裝自己的思維，用工具放大自己的能力——這就是 AI 程式設計的全部秘密。

\booksection{參考文獻}

\begin{enumerate}
  \item Andrej Karpathy. (2025). Vibe Coding. Social media post introducing the concept of vibe-driven programming with AI.
  \item Andrej Karpathy. (2025). CLAUDE.md - Project guidelines for AI-assisted development. GitHub repository, 70+K stars.
  \item Anthropic. (2026). Claude Code: Terminal-native AI programming. Official documentation and best practices.
  \item Cursor Team. (2026). Cursor 3: The next generation of AI-first IDE. Release notes and documentation.
  \item OpenAI. (2026). Codex CLI: Bringing AI to the terminal. Technical documentation.
  \item Model Context Protocol (MCP) Specification. (2025-2026). Open protocol for AI-tool integration.
  \item Vibe Coding Community. (2026). Community-curated skills and best practices for AI-assisted programming.
  \item Spec-Driven Development (SDD) methodology. (2026). Emerging practices in specification-first AI programming.
\end{enumerate}
